
\documentclass[11pt,a4paper]{report}

% ==================================================
% 1. PACKAGES
% ==================================================

\usepackage[utf8]{inputenc}
\usepackage[T1]{fontenc}
\usepackage[french]{babel}
\usepackage{lmodern}
\usepackage{geometry}
\usepackage{graphicx}
\usepackage{float}
\usepackage{array}
\usepackage{booktabs}
\usepackage{longtable}
\usepackage{tabularx}
\usepackage{enumitem}
\usepackage{hyperref}
\usepackage{xcolor}
\usepackage{fancyhdr}
\usepackage{titlesec}
\usepackage{setspace}
\usepackage{caption}
\usepackage{placeins}

% ==================================================
% 2. CONFIGURATION DU DOCUMENT
% ==================================================

\geometry{
    left=2.5cm,
    right=2.5cm,
    top=2.2cm,
    bottom=2.2cm
}

\setstretch{1.15}
\setlength{\parindent}{0.5cm}
\setlength{\parskip}{4pt}

\hypersetup{
    colorlinks=true,
    linkcolor=black,
    urlcolor=blue,
    citecolor=black
}

\pagestyle{fancy}
\fancyhf{}
\fancyhead[L]{CELEBRITY CITY}
\fancyhead[R]{Cahier des charges}
\fancyfoot[C]{\thepage}

% ==================================================
% 3. COMMANDE POUR LES DIAGRAMMES
% ==================================================

% Utilisation :
% \diagramme{fichier.png}{Nom du diagramme}{Description}

\newcommand{\diagramme}[3]{
    \begin{figure}[H]
        \centering
        \IfFileExists{#1}{
            \includegraphics[
                width=0.90\textwidth,
                height=0.36\textheight,
                keepaspectratio
            ]{#1}
        }{
            \fbox{
                \parbox[c][4cm][c]{0.85\textwidth}{
                    \centering
                    Diagramme à insérer : #2\\
                    Fichier attendu : #1
                }
            }
        }
        \caption{#2}
        \label{fig:#2}
    \end{figure}
    \textit{Description : #3}
}

% ==================================================
% 4. PAGE DE GARDE
% ==================================================

\begin{document}

\begin{titlepage}
    \begin{center}

        \textbf{RÉPUBLIQUE DU CAMEROUN}\\
        Paix - Travail - Patrie

        \vspace{1cm}

        \textbf{[NOM DE L'ÉTABLISSEMENT]}\\
        \textbf{[DÉPARTEMENT / FILIÈRE]}

        \vspace{1.5cm}

        % Logo de l'établissement
        % \includegraphics[width=3cm]{images/logo.png}

        \vspace{1cm}

        {\Huge\bfseries CAHIER DES CHARGES}

        \vspace{1cm}

        {\LARGE\bfseries CELEBRITY CITY}

        \vspace{0.5cm}

        {\large
        Système de signalement citoyen
        et de gestion de la collecte des déchets
        }

        \vspace{1.5cm}

        \textbf{Présenté par :}

        \vspace{0.3cm}

        [Noms des membres de l'équipe]

        \vfill

        \textbf{Encadrant :} [Nom de l'enseignant]

        \vspace{0.5cm}

        Année académique : [2026--2027]

    \end{center}
\end{titlepage}

% ==================================================
% 5. PAGES PRÉLIMINAIRES
% ==================================================

\pagenumbering{roman}

\chapter*{Avant-propos}
\addcontentsline{toc}{chapter}{Avant-propos}

% Présenter brièvement le cadre académique du projet.

\chapter*{Résumé}
\addcontentsline{toc}{chapter}{Résumé}

% Résumer le problème, la solution et les objectifs du projet.

\chapter*{Liste des sigles et abréviations}
\addcontentsline{toc}{chapter}{Liste des sigles}

% API, GPS, UML, SGBD, etc.

\tableofcontents
\listoffigures
\listoftables

\clearpage
\pagenumbering{arabic}

% ==================================================
% CHAPITRE 1 : PRÉSENTATION GÉNÉRALE
% ==================================================

\chapter{Présentation générale du projet}

\section{Contexte du projet}

\section{Justification du projet}

\section{Problématique}

\section{Présentation de Celebrity City}

\section{Objectif général}

\section{Objectifs spécifiques}

\section{Résultats attendus}

% ==================================================
% CHAPITRE 2 : ÉTUDE DE L'EXISTANT
% ==================================================

\chapter{Étude de l'existant}

\section{Présentation de la situation actuelle}

\section{Analyse des méthodes de signalement existantes}

\section{Analyse des méthodes de collecte des déchets}

\section{Difficultés et insuffisances constatées}

\section{Solution proposée}

% ==================================================
% CHAPITRE 3 : PÉRIMÈTRE DU PROJET
% ==================================================

\chapter{Délimitation et périmètre du projet}

\section{Périmètre fonctionnel}

\section{Fonctionnalités retenues}

\section{Fonctionnalités exclues}

\section{Périmètre géographique}

\section{Limites du projet}

% ==================================================
% CHAPITRE 4 : ACTEURS ET BESOINS
% ==================================================

\chapter{Identification des acteurs et analyse des besoins}

\section{Identification des acteurs}

\subsection{Citoyen}

\subsection{Agent municipal}

\subsection{Responsable municipal}

\subsection{Responsable de collecte}

\subsection{Administrateur}

\section{Besoins fonctionnels par acteur}

\section{Besoins et attentes des utilisateurs}

% DIAGRAMME :
% Diagramme de contexte / acteurs du système

\diagramme{diagrams/01-contexte.png}
{Diagramme de contexte}
{Représentation des acteurs externes et de leurs interactions avec Celebrity City.}

% ==================================================
% CHAPITRE 5 : SPÉCIFICATIONS FONCTIONNELLES
% ==================================================

\chapter{Spécifications fonctionnelles}

\section{Module d'authentification et de gestion des comptes}

\subsection{Inscription et connexion}

\subsection{Gestion des rôles et permissions}

\section{Module de signalement citoyen}

\subsection{Création d'un signalement}

\subsection{Ajout de photographies}

\subsection{Géolocalisation automatique}

\subsection{Catégorisation des problèmes}

\subsection{Consultation et suivi des signalements}

\section{Module de gestion des signalements}

\subsection{Consultation des signalements}

\subsection{Vérification et validation}

\subsection{Affectation aux équipes municipales}

\subsection{Gestion des statuts}

\subsection{Notifications}

\section{Module de gestion de la collecte}

\subsection{Gestion des zones de collecte}

\subsection{Gestion des camions et des agents}

\subsection{Planification des collectes}

\subsection{Suivi des opérations}

\section{Module cartographique}

\subsection{Affichage des signalements}

\subsection{Localisation des poubelles publiques}

\subsection{Visualisation des camions}

\subsection{Cartographie des zones à problèmes}

\section{Tableau de bord municipal}

\subsection{Statistiques des signalements}

\subsection{Suivi des interventions}

\subsection{Indicateurs de collecte}

% ==================================================
% CHAPITRE 6 : SPÉCIFICATIONS NON FONCTIONNELLES
% ==================================================

\chapter{Spécifications non fonctionnelles}

\section{Sécurité}

\section{Performance}

\section{Disponibilité}

\section{Ergonomie et accessibilité}

\section{Compatibilité multiplateforme}

\section{Maintenabilité et évolutivité}

\section{Gestion des connexions instables}

% ==================================================
% CHAPITRE 7 : MODÉLISATION UML
% ==================================================

\chapter{Modélisation et conception UML}

\section{Présentation de UML}

\section{Diagramme de cas d'utilisation général}

\diagramme{diagrams/02-use-case-global.png}
{Diagramme de cas d'utilisation général}
{Présentation des acteurs et de l'ensemble des fonctionnalités accessibles dans le système.}

\section{Diagrammes de cas d'utilisation détaillés}

\subsection{Cas d'utilisation du citoyen}

\diagramme{diagrams/03-use-case-citoyen.png}
{Diagramme de cas d'utilisation du citoyen}
{Inscription, connexion, création et suivi des signalements.}

\subsection{Cas d'utilisation du personnel municipal}

\diagramme{diagrams/04-use-case-municipalite.png}
{Diagramme de cas d'utilisation municipal}
{Vérification, affectation et traitement des signalements.}

\subsection{Cas d'utilisation de la collecte}

\diagramme{diagrams/05-use-case-collecte.png}
{Diagramme de cas d'utilisation de la collecte}
{Planification et suivi des opérations de collecte.}

\section{Diagrammes d'activités}

\subsection{Processus de création d'un signalement}

\diagramme{diagrams/06-activite-signalement.png}
{Diagramme d'activités du signalement}
{Étapes allant de la création du signalement à son enregistrement.}

\subsection{Processus de traitement d'un signalement}

\diagramme{diagrams/07-activite-traitement.png}
{Diagramme d'activités du traitement}
{Circuit de vérification, d'affectation et de résolution d'un signalement.}

\section{Diagrammes de séquence}

\subsection{Séquence de création d'un signalement}

\diagramme{diagrams/08-sequence-signalement.png}
{Diagramme de séquence du signalement}
{Échanges entre le citoyen, l'application mobile, l'API et la base de données.}

\subsection{Séquence d'affectation d'une intervention}

\diagramme{diagrams/09-sequence-affectation.png}
{Diagramme de séquence de l'affectation}
{Échanges entre le responsable municipal, le backend et l'agent affecté.}

\section{Diagramme de classes}

\diagramme{diagrams/10-classes.png}
{Diagramme de classes}
{Présentation des classes métier, attributs, méthodes et relations du système.}

\section{Diagramme de composants}

\diagramme{diagrams/11-composants.png}
{Diagramme de composants}
{Organisation des composants Laravel, React, Flutter et des services externes.}

\section{Diagramme de déploiement}

\diagramme{diagrams/12-deploiement.png}
{Diagramme de déploiement}
{Représentation de l'environnement d'exécution et des communications entre les applications et les serveurs.}

% ==================================================
% CHAPITRE 8 : BASE DE DONNÉES
% ==================================================

\chapter{Conception de la base de données}

\section{Choix du système de gestion de base de données}

\section{Identification des entités}

\section{Description des entités et attributs}

\section{Relations entre les entités}

\section{Règles de gestion}

\section{Modèle conceptuel de données (MCD)}

\diagramme{diagrams/13-mcd.png}
{Modèle conceptuel de données}
{Présentation des entités, associations, cardinalités et règles de gestion.}

\section{Modèle logique de données (MLD)}

% ==================================================
% CHAPITRE 9 : ARCHITECTURE TECHNIQUE
% ==================================================

\chapter{Architecture technique du système}

\section{Architecture générale}

\section{Backend : Laravel}

\section{Frontend web : React}

\section{Application mobile : Flutter}

\section{Communication par API REST}

\section{Communication temps réel par WebSocket}

\section{Services cartographiques et géolocalisation}

\section{Choix des technologies et justification}

% ==================================================
% CHAPITRE 10 : INTERFACES UTILISATEUR
% ==================================================

\chapter{Conception des interfaces utilisateur}

\section{Interface d'authentification}

\section{Interface citoyenne}

\section{Interface de création d'un signalement}

\section{Interface de suivi des signalements}

\section{Interface d'administration municipale}

\section{Interface de gestion de la collecte}

\section{Tableau de bord et cartographie}

% Emplacements des maquettes :
% images/maquette-connexion.png
% images/maquette-citoyen.png
% images/maquette-signalement.png
% images/maquette-carte.png
% images/maquette-dashboard.png

% ==================================================
% CHAPITRE 11 : ORGANISATION DU PROJET
% ==================================================

\chapter{Organisation et planification du projet}

\section{Méthodologie de développement}

\section{Découpage en phases}

\section{Répartition des responsabilités}

\section{Outils de collaboration et de gestion}

\section{Calendrier prévisionnel}

\section{Gestion des versions avec Git et GitHub}

% ==================================================
% CHAPITRE 12 : TESTS ET VALIDATION
% ==================================================

\chapter{Stratégie de tests et validation}

\section{Tests unitaires}

\section{Tests d'intégration}

\section{Tests fonctionnels}

\section{Tests de sécurité}

\section{Tests de performance}

\section{Critères d'acceptation}

% ==================================================
% CHAPITRE 13 : CONTRAINTES ET RISQUES
% ==================================================

\chapter{Contraintes et gestion des risques}

\section{Contraintes techniques}

\section{Contraintes matérielles et financières}

\section{Risques liés à la géolocalisation}

\section{Risques liés à la connectivité}

\section{Mesures de prévention}

% ==================================================
% CHAPITRE 14 : CONCLUSION
% ==================================================

\chapter{Conclusion générale}

% Résumer le projet, sa pertinence et les résultats attendus.

% ==================================================
% RÉFÉRENCES ET ANNEXES
% ==================================================

\begin{thebibliography}{9}

\bibitem{uml}
Référence à la documentation UML.

\bibitem{laravel}
Documentation officielle Laravel,
\url{https://laravel.com/docs}

\bibitem{react}
Documentation officielle React,
\url{https://react.dev}

\bibitem{flutter}
Documentation officielle Flutter,
\url{https://docs.flutter.dev}

\bibitem{osm}
Documentation OpenStreetMap,
\url{https://www.openstreetmap.org}

\end{thebibliography}

\appendix

\chapter{Glossaire}

\chapter{Règles de gestion complémentaires}

\chapter{Documents complémentaires}

\end{document}
